\chapter{Algoritmos e estruturas de dados aplicados}

\section{list / array}

Em Python, \code{list} e um array dinamico. No projeto, ela aparece como batch:

\begin{lstlisting}[style=devbutterpython]
batch: list[HealthRecord] = []
batch.append(record)
\end{lstlisting}

\code{append} e O(1) medio. Acesso por indice tambem e O(1).

\section{dict / hash map}

\code{dict} e uma tabela hash. Usamos para metadata e contadores.

\begin{lstlisting}[style=devbutterpython]
metadata: dict[str, str] = {}
metadata[key] = value
\end{lstlisting}

Lookup medio:

\[
O(1)
\]

\section{set / hash set}

\code{set} e usado para membership rapido.

\begin{lstlisting}[style=devbutterpython]
allowed_types = {
    "HKQuantityTypeIdentifierHeartRate",
    "HKQuantityTypeIdentifierStepCount",
}

if record.type in allowed_types:
    process(record)
\end{lstlisting}

Checar se um item esta no set e O(1) medio.

\section{queue}

Fila FIFO usada em producer/consumer.

\begin{lstlisting}[style=devbutterpython]
record_queue.put(record)
record = record_queue.get()
\end{lstlisting}

\section{tree}

XML e uma arvore:

\begin{verbatim}
Workout
 ├── MetadataEntry
 ├── WorkoutStatistics
 └── WorkoutEvent
\end{verbatim}

Cada tag e um no; tags internas sao filhos.

\section{recursion}

Recursao pode percorrer arvores pequenas:

\begin{lstlisting}[style=devbutterpython]
def walk(element):
    print(element.tag)

    for child in element:
        walk(child)
\end{lstlisting}

Mas evite usar recursao carregando a arvore inteira do XML gigante. Para o export real, prefira streaming.

\section{sorting}

Ordenar por data e comum em series temporais. Ordenacao geralmente custa:

\[
O(n \log n)
\]

No banco:

\begin{lstlisting}[style=devbuttersql]
SELECT *
FROM health_records
ORDER BY start_date;
\end{lstlisting}

\section{binary search}

Busca binaria encontra um item em lista ordenada em O(log n). Em Python:

\begin{lstlisting}[style=devbutterpython]
from bisect import bisect_left

positions = [1, 3, 5, 7, 9]
index = bisect_left(positions, 6)
\end{lstlisting}

No projeto, a ideia aparece ao buscar registros em series temporais ordenadas. No banco, indices B-tree usam uma logica parecida de busca eficiente.

\section{two pointers}

Dois ponteiros servem para cruzar duas listas ordenadas sem fazer O(n²).

Exemplo: achar batimentos dentro de intervalos de sono.

\begin{verbatim}
sono:       [inicio, fim]
heart rate: timestamp ordenado
\end{verbatim}

Em vez de comparar todo sono com todo batimento, voce avanca ponteiros conforme as datas.

\section{sliding window}

Sliding window e uma janela movel. Exemplo: media movel de 7 dias.

Em SQL:

\begin{lstlisting}[style=devbuttersql]
SELECT
    day,
    steps,
    AVG(steps) OVER (
        ORDER BY day
        ROWS BETWEEN 6 PRECEDING AND CURRENT ROW
    ) AS steps_7d_avg
FROM daily_summary
ORDER BY day;
\end{lstlisting}

Isso e muito util para tendencias de passos, sono e batimento.

\section{O(n²): o que evitar}

Evite loops aninhados sem necessidade:

\begin{lstlisting}[style=devbutterpython]
for sleep_interval in sleep_intervals:
    for heart_rate in heart_rates:
        compare(sleep_interval, heart_rate)
\end{lstlisting}

Se cada lista tem muitos itens, isso vira O(n²). Busque alternativas com ordenacao, indices, SQL, two pointers ou janelas.
